\input{../tex-inputs/latex-leseansicht-vorspann}

\section[Arthur Schnitzler an Gustav Schwarzkopf, {[}zwischen 4. und 10. 10. 1894?{]}]{L04211 Arthur Schnitzler an Gustav Schwarzkopf, {[}zwischen 4. und 10. 10. 1894?{]}}
\nopagebreak\mylabel{L04211v}
\rehead{ }\normalsize\beginnumbering\briefempfaengerindex{Schwarzkopf, Gustav@\textsc{Schwarzkopf, Gustav}!zzzSchnitzler, Arthur@\emph{von Arthur Schnitzler}!1894-10-101@{{[}zwischen 4. und 10. 10. 1894?{]}}|(be}
\toendnotes[C]{\smallbreak\pagebreak[2]}
\correspDesc{Versand  durch Arthur Schnitzler im Zeitraum [zwischen 4. und
                  10. 10. 1894?] in Frankgasse 1
\newline{}Erhalt  durch Gustav Schwarzkopf in Tiefer Graben 23}\toendnotes[C]{\smallbreak}
\Standort{CUL, Schnitzler, B 96.}
\physDesc{Briefkarte, 385 Zeichen
\newline{}Handschrift: Bleistift, deutsche Kurrent}\toendnotes[C]{\smallbreak}
\pstart
           \noindent{}{\pb}Lieber Freund, ich wiederhole meine Bitte um Schonungsloſigkeit\pwindex{Schnitzler, Arthur 15. 5. 1862 Wien – 21. 10. 1931 ebd.@\textsc{Schnitzler, Arthur} (15. 5. 1862 Wien – 21. 10. 1931 ebd.)!kleine Komödie@\strich\emph{Die kleine Komödie}|pwv}. Vergeſſen Sie nicht,
               dß abſolut kein zwingender Grund für die Veröffentlichg von was ſchlechtem oder
               mittelmäßigem vorliegt. We{\geminationn} Sie zu Ende geleſen haben,
               bevor wir {\pb}einander wo
               begegnen, ſo ſchreiben Sie mir eine Zeile und ich komme zu Ihrer Eſſenszeit hinauf,
               mir \label{K_L04211-1v}\edtext{\textsc{Manuscript\pwindex{Schnitzler, Arthur 15. 5. 1862 Wien – 21. 10. 1931 ebd.@\textsc{Schnitzler, Arthur} (15. 5. 1862 Wien – 21. 10. 1931 ebd.)!kleine Komödie@\strich\emph{Die kleine Komödie}|pwv}} und Urtheil}{\lemma{\textnormal{\emph{Manuscript und Urtheil}}}\Cendnote{\textnormal{Das Korrespondenzstück
                  ist undatiert. Es fügt sich aber inhaltlich an das Schreiben vom XXXX Auszeichnungsfehler: Dokument L04105 nicht gefunden. Am 10. 10. 1894 erhielt
                     Schnitzler das gewünschte Urteil, was
                  zeitlich den Punkt markiert, an dem der Brief bereits übermittelt gewesen sein
                  muss. Für die generelle Einordnung in die zeitliche Epoche spricht die Verwendung
                  von »Lieber Freund« als Anrede, die nur zwischen XXXX Auszeichnungsfehler: Dokument L04105 nicht gefunden und XXXX Auszeichnungsfehler: Dokument L04116 nicht gefunden in der
                  Korrespondenz zum Einsatz kommt.}}}\label{K_L04211-1} abzuholen.\pend
           
\pstart
           Herzlich grüßend Ihr{\\[\baselineskip]}\spacefill\mbox{ArthSchn}\pend
           \leftskip=0em{}\selectlanguage{ngerman}\endnumbering\briefempfaengerindex{Schwarzkopf, Gustav@\textsc{Schwarzkopf, Gustav}!zzzSchnitzler, Arthur@\emph{von Arthur Schnitzler}!1894-10-041@{{[}zwischen 4. und 10. 10. 1894?{]}}|)be}\mylabel{L04211h}
\begin{anhang}
\end{anhang}\newcommand{\dateiname}{L04211}\newcommand{\titel}{Arthur Schnitzler an Gustav Schwarzkopf, [zwischen 4. und 10. 10. 1894?]}\newcommand{\editorInnen}{Herausgegeben von Jahnke, SelmaMüller, Martin Anton}\input{../tex-inputs/latex-leseansicht-abspann}
